\subsection{Alignment denominators, grouping audit, and deterministic folds}
\label{sec:fold_details}
Observed sequences are assembled separately for each chain. Biopython local pairwise alignment uses BLOSUM62, gap opening $-10$, gap extension $-0.5$, and the first optimal alignment returned by the fixed implementation. Identity is the number of exactly matching non-X residue pairs divided by the number of aligned nongap residue pairs. Coverage for each chain is the same nongap pair count divided by that chain's observed length. Thus X contributes to the denominator when aligned but does not count as an identity match. A chain pair qualifies at identity $\geq0.30$, reciprocal coverage $\geq0.80$, and at least 50 aligned nongap residue pairs.

Any qualifying chain pair links the corresponding whole protein entries. Exact complete chain-sequence multiset duplicates are also linked irrespective of size; unknown positions use their explicit three-letter residue identities for this duplicate check, preventing distinct X-coded components from becoming artificial matches. Identical small auxiliary chains alone do not activate the exception. Connected components define the final grouping. This is conservative with respect to proteins sharing a qualifying chain and allows chains within a multimer to keep their common structural context. It is not a claim that a 30\% threshold alone establishes or excludes evolutionary homology. Entries containing only chains shorter than 50 residues rely on the complete-sequence duplicate rule and receive no assurance against more distant sequence similarity.

The initial grouping rule used coverage of only the shorter chain and had no minimum aligned length. On the initial standard-residue audit, it produced 3,434 edges and a single component containing all 364 eligible entries, driven by transitive bridges through very short auxiliary fragments. Before prediction outcomes were accessed, the rule was amended to the reciprocal-coverage and minimum-length definition above. After the final chemical-identity audit, the graph has 16 edges and 349 components: 336 singletons, 11 pairs, and two triplets. The final chain audit records 565 observed chains, of which 108 are shorter than 50 residues; 47 entries have only such short chains, and 119 entries have multiple chains. Both the original graph and the amendments are retained as protocol evidence. These audited counts describe the grouping process, not a measured improvement in prediction.

Cluster components are sorted by decreasing total residue count, with a seeded SHA-256 tie-breaking key. In order, each is assigned to the currently least-loaded of five folds, breaking load ties by numeric fold identifier. Ordinary protein-only folds use the same rule on individual proteins. The assignment is deterministic and label-free. A fold manifest is shared across feature families, so a comparison cannot improve by drawing a more favorable partition. The bootstrap resamples the resulting whole clusters 2,000 times, conditional on the held-out predictions, and recomputes the macro-protein estimand in every resample.

Pearson and Spearman correlations are marked undefined if fewer than three paired residues are available or the target or prediction has standard deviation at most $10^{-12}$. Such entries remain in prediction and RMSE tables but contribute no finite correlation to its macro mean. Each metric summary records its number of valid proteins; paired correlation differences use only proteins with finite values for both models. The bootstrap first retains finite values for the specified estimand, then samples whole remaining clusters with replacement and reports the 2.5th and 97.5th percentiles of the resulting protein means. Invalid predictions or unmatched protein cohorts cause an error instead of silent row deletion.
